\documentclass[10pt,a4paper]{amsart}
\usepackage[margin=19mm]{geometry}
\usepackage{amsmath,amssymb,mathtools}
\usepackage[scr=rsfs,cal=euler]{mathalfa}
\usepackage{xfrac}
\usepackage[unicode,hidelinks,bookmarks=false]{hyperref}
\newcommand{\ie}{i.e.\ }
\newcommand{\cf}{cf.\ }
\newcommand{\resp}{respectively\ }
\newcommand{\Subsection}{\subsection}
\providecommand{\nobreakdash}{\nobreak\hbox{-}}
\setlength{\emergencystretch}{2em}
\multlinegap0pt
\usepackage{bm}
\usepackage{bbm}
\usepackage{dsfont}
\usepackage{xspace}
\usepackage{cancel}
\usepackage{mathtools}
\usepackage{amsthm}
\makeatletter
\@ifundefined{teorema}{\newtheorem{teorema}{Teorema}}{}
\@ifundefined{lema}{\newtheorem{lema}[teorema]{Lemma}}{}
\makeatother
\newcommand{\IdG}{\ensuremath{\textnormal{Id}^{(G,\ast)}}}
\newcommand{\spam}{\ensuremath{\textnormal{span}}}
\title[Algebras with additional structures and multiplicities bound]{Algebras with additional structures and multiplicities bounded by a constant}
\author{R. B. dos Santos, A. C Vieira, R. F. D. N. Vieira}
\date{Source: arXiv:2510.03583v1 (CC BY 4.0)}
\begin{document}
\maketitle

\noindent\emph{Source and licence.} Algebras with additional structures and multiplicities bounded by a constant. Version arXiv:2510.03583v1 (CC BY 4.0), distributed under the Creative Commons Attribution 4.0 International licence, as recorded in the arXiv licence metadata for this exact version and shown on the abstract page https://arxiv.org/abs/2510.03583v1. Full rights notice: licenses/arxiv-2510.03583-CC-BY-4.0.txt.\par\smallskip

\noindent\textit{Editorial scope.} Counted unit: the Lema whose source label is \texttt{lema2F} in the article (source0.tex lines 928--938, source label \texttt{lema2F}), delivered together with the article's own complete original proof of it (source0.tex lines 939--1003, one \texttt{proof} environment of 65 lines). No mathematics is written, repaired or re-derived: statement and proof are the article's own text line for line. Required context retained uncounted: lema1F (lines 806--813); lemmanovo (lines 827--828); ref0.1 (lines 844--854); eq-1 (lines 845--847); eq-2 (lines 848--850); ref0.2 (lines 888--891); ref0.3 (lines 901--903). Every definition, display and label of those lines is retained unchanged. Omitted from this package and not redistributed: 1--805, 814--826, 829--843, 851--887, 892--900, 904--927, 1004--1227. Nothing inside a retained excerpt is omitted or altered: every retained body is delivered exactly as the article's lines are written, so no omission receipt of any kind accompanies this package. Citations: the retained excerpts carry no citation command at all, so no bibliography entry of the article is transcribed and no reference is redistributed. Reference adaptation: none was needed and none was made; every reference inside the delivered unit resolves inside it, so no unresolved reference or citation is produced. Printed slips kept literal: no printed word, symbol, hypothesis or number of the counted statement or of its proof was altered. Layout and class: the article's own class is \texttt{amsart} with packages including fullpage, array, mathrsfs, hyperref, alltt, graphicx, graphpap, color, ulem, inputenc, amsfonts, amssymb, verbatim; the wrapper is the lane's pinned \texttt{amsart} editorial wrapper at 10pt on A4, which restates the theorem-like environments the retained text needs and adds the article's own macro definitions that the retained text needs (\texttt{\textbackslash IdG}, \texttt{\textbackslash spam}), with extra wrapper packages (bm, bbm, dsfont, xspace, cancel, mathtools). The delivered numbering is the unit's own. Assets: no retained span contains a figure environment, a graphics-inclusion command, a PDF-page inclusion, a table or a TikZ picture; no image file is copied into this package, the package declares no asset dependency and no image file is redistributed with it. Licence: version arXiv:2510.03583v1 of this article is distributed under the Creative Commons Attribution 4.0 International licence, as recorded in the arXiv licence metadata for this exact version and shown on the abstract page https://arxiv.org/abs/2510.03583v1. A full rights notice is delivered at licenses/arxiv-2510.03583-CC-BY-4.0.txt. Characters: every retained excerpt is pure ASCII, so the delivered engine, XeLaTeX with its default Latin Modern text font, reports no missing character.
\begin{lema}\label{lema1F} 
 Let $g\in G\backslash \{1\}$ and assume that $A$ satisfies at least one of the identities
	\begin{center}
		$y_{1,g^2}y_{2,g}\equiv 0$ \;  or \;$z_{1,g^2}y_{2,g}\equiv 0.$
	\end{center}
	Then $m_{\langle \lambda \rangle}\le 1$, where $\langle \lambda \rangle$ is a multipartition of type $(\mu_{g^+})\vdash (n_{g^+})$.

\end{lema} 

\begin{lema}\label{lemmanovo} Let $g\in G\backslash \{1\}$ and suppose that $A$ satisfies the identity $y_{1,g}y_{3,g}y_{2,g}+y_{2,g}y_{3,g}y_{1,g}\equiv 0$. Then we have $m_{\langle \lambda \rangle}\le 1$, for all multipartition $\langle \lambda \rangle \vdash (3_{g^+})$.  
		\end{lema} 

\begin{lema}\label{ref0.1} Let $g\in G\backslash \{1\}$ and suppose that $A$ satisfies the following identities 
\begin{equation}\label{eq-1}
    y_{1,g}y_{3,g}y_{2,g}+y_{2,g}y_{3,g}y_{1,g}\equiv 0\;\;\;
\end{equation}
\begin{equation}\label{eq-2}
\;\;\;y_{1,g}y_{2,g}y_{4,g}y_{3,g}+ y_{2,g}y_{4,g}y_{3,g} y_{1,g} \equiv 0
.\end{equation} 
Then we have $m_{\langle \lambda \rangle}\le 1$, for all multipartition $\langle \lambda \rangle \vdash (4_{g^+})$.

     
		\end{lema}

 \begin{lema}\label{ref0.2} 
Assume that $A$ satisfies the identities (\ref{eq-1}) and (\ref{eq-2}). Then, for $n>4,$ we have 
  	$P_{(n_{g^+})}(A)= \spam_F\{y_{1,g}y_{2,g}\cdots y_{n,g}\}$. Consequently, $m_{\langle \lambda \rangle}\le 1$, for all multipartition $\langle \lambda \rangle \vdash (n_{g^+})$ with $n> 4$.
	\end{lema}

 \begin{lema}{\label{ref0.3}} 
Assume that $A$ satisfies the identity $y_{1,g}y_{3,g}y_{2,g}-y_{2,g}y_{1,g}y_{3,g}\equiv 0.$ Then, for $n\ge 3$ we have $m_{\langle \lambda \rangle}\le 1$, for all multipartition $\langle \lambda \rangle \vdash (n_{g^+}).$
	\end{lema} 

\begin{lema}\label{lema2F} 
 Let $g\in G\backslash \{1\}$  and assume that $A$ satisfies at least one identity in each one of the following items
	\begin{itemize}
		\item[1)]$y_{1,g^2}y_{2,g}+\gamma_{i} y_{2,g}y_{1,g^2}\equiv 0;$
		
		\item[2)]$z_{1,g^2}y_{2,g}+\gamma_{l} y_{2,g}z_{1,g^2} \equiv 0;$
		
		\item[3)] $y_{1,g}z_{2,g^3}+\gamma_{m} z_{2,g^3}y_{1,g}\equiv 0;$
		
	\end{itemize} where $\gamma_{i}, \gamma_{l},$ $\gamma_{m} \in\{ 0,1,-1\}.$ Then, for $n\geq 3,$ we have $m_{\langle \lambda \rangle}\le 1$, for all multipartition $\langle \lambda \rangle=(\mu_{g^+})$ of  $(n_{g^+})$.
\end{lema} 

\begin{proof}

    Initially, note that if $\gamma_{i} = 0$ ($\gamma_l=0$, respec.) then $A$ satisfies the identity $y_{1,g^2}y_{2,g}\equiv 0$ ($z_{1,g^2}y_{2,g}\equiv 0$, respec.). Therefore, the result follows from Lemma \ref{lema1F}. Next we study the other cases.
    
    Suppose that $A$ satisfies the identity of item $1)$ with $\gamma_{i}=1$. We then move on to analyze the identities in item $2)$. 
    If $\gamma_{l}=0,$ we have  $y_{1,g^2}y_{2,g}+y_{2,g}y_{1,g^2} \equiv 0 \; \mbox{and}\;z_{1,g^2}y_{2,g}  \equiv 0. $ In this case, the result follows from Lemma \ref{lema1F}. 

   	Now if $\gamma_{l}=1,$ we have  
$y_{1,g^2}y_{2,g}+y_{2,g}y_{1,g^2}\equiv 0 \;\;\mbox{and}\;\ z_{1,g^2}y_{2,g}+y_{2,g}z_{1,g^2}\equiv 0.$  
Since $y_{1,g}\circ y_{3,g}$ and $[y_{1,g},y_{3,g}]$ are, respectively, symmetric and skew $(G,*)$-polynomials of homogeneous degree $g^2$, we have
$$(y_{1,g}\circ y_{3,g})y_{2,g}+y_{2,g}(y_{1,g}\circ y_{3,g})\equiv 0 \;\;\;\mbox{and}\;\;\; [y_{1,g},y_{3,g}]y_{2,g}+y_{2,g}[y_{1,g},y_{3,g}]\equiv 0.$$    

 Adding the two identities given above, we get 
$y_{1,g}y_{3,g}y_{2,g}+y_{2,g}y_{1,g}y_{3,g}\equiv 0.$  
 Thus,
$$ 
    \begin{array}{ccl}y_{1,g}y_{3,g}y_{2,g}& \equiv& -y_{2,g}y_{1,g}y_{3,g}\\ & \equiv & y_{3,g}y_{2,g}y_{1,g} \\ & \equiv& -y_{1,g}y_{3,g}y_{2,g}\end{array}$$
	and this implies that  $y_{1,g}y_{3,g}y_{2,g}\in \IdG(A)$.
	Therefore, for $n \ge 3$ we have that $P_{(n_{g^+})} \subseteq \IdG(A),$ and so $\dim_FP_{(n_{g^+})}(A)=0$. Consequently, we conclude that $m_{\langle \lambda \rangle}= 0$ in this case, for every multipartition $\langle \lambda \rangle$ of type $(\mu_{g^+})\vdash (n_{g^+})$, and so the lemma follows. 
     

	Now we consider the case $\gamma_{i}=1$ and $\gamma_{l}=-1$. In this case, we have that $A$ satisfies the identities 	$	y_{1,g^2}y_{2,g}+y_{2,g}y_{1,g^2}\equiv 0$ and 
$	z_{1,g^2}y_{2,g}-y_{2,g}z_{1,g^2}\equiv 0. $
  Similarly to the previous case, using the polynomials $y_{1,g}\circ y_{3,g}$ and $[y_{1,g},y_{3,g}]$, as a consequence,  we obtain the identity 
    	\begin{equation} \label{eq7F}
y_{1,g}y_{3,g}y_{2,g}+y_{2,g}y_{3,g}y_{1,g} \equiv 0.
	\end{equation}	
    
    If $n=3$, according to Lema \ref{lemmanovo} we get $m_{\langle \lambda \rangle}\le 1$, for every multipartition $\langle \lambda \rangle$ of type $(\mu_{g^+})\vdash (3_{g^+})$.  Thus, we will treat the case where $n\geq 4$ and consider the identities that appear in item $3).$  
    
    Suppose that $\gamma_{m}=0,$ i.e. $A$ satisfies the identity 
\begin{equation}\label{eq8F} y_{1,g}z_{2,g^3} \equiv 0.
	\end{equation} 

    Define the $(G,*)$-polynomial $f=y_{2,g}y_{4,g}y_{3,g}-y_{3,g}y_{4,g}y_{2,g}$ and note that $f$ is skew of degree $g^3$. Using the identity (\ref{eq8F}) we get $y_{1,g}f\equiv 0$ on $A$ and so $y_{1,g}y_{2,g}y_{4,g}y_{3,g}-y_{1,g}y_{3,g}y_{4,g}y_{2,g}\in \IdG(A)$. Again using the identity (\ref{eq7F}) we get 
$$y_{1,g}y_{2,g}y_{4,g}y_{3,g}\equiv 0. $$ 

Therefore, for $n\geq 4,$ we have $P_{(n_{g^+})}\subset \IdG(A)$ and thus $m_{\langle \lambda \rangle}= 0$, for every multipartition $\langle \lambda \rangle $ of type $(\mu_{g^+})\vdash (n_{g^+})$. 

 Next consider the situation where $\gamma_{i}=1,$ $\gamma_{l}=-1$ and $\gamma_{m}=1$. In this case, we have the identity 
	$$y_{1,g}z_{2,g^3} + z_{2,g^3}y_{1,g}\equiv 0.$$ 
    
    Considering the endomorphism that takes the variable $z_{2,g^3}$ in the polynomial $f$ defined above and using the identity (\ref{eq7F}),  as a consequence of the above identity, we obtain  
\begin{equation}\label{eq9F} y_{1,g}y_{2,g}y_{4,g}y_{3,g}+y_{2,g}y_{4,g}y_{3,g}y_{1,g}\equiv 0.
	\end{equation} 

Thus,  by Lemmas \ref{ref0.1} and \ref{ref0.2} it follows that $m_{\langle \lambda \rangle}\le 1$, for every $\langle \lambda \rangle \vdash (n_{g^+})$ with $n\geq 4.$
	
In the situation where $\gamma_{i}=1,$ $\gamma_{l}=-1$ and $\gamma_{m}=-1,$ 
	the proof is analogous to the last case. 
    Then, it remains to consider the case $\gamma_i=-1.$ 
     If $\gamma_{l}=1$ we have that 
	$y_{1,g^2}y_{2,g}-y_{2,g}y_{1,g^2}\equiv 0$
		and $z_{1,g^2}y_{2,g}+y_{2,g}z_{1,g^2}\equiv 0.$
In this case, we proceed analogously to the case  $\gamma_{i}=1$ and $\gamma_{l}=-1.$ 

Finally, we consider the case with $\gamma_{i}=-1$ and $\gamma_{l}=-1,$ i.e. we have the following identities
	$$y_{1,g^2}y_{2,g}-y_{2,g}y_{1,g^2}\equiv 0,$$
	$$ z_{1,g^2}y_{2,g}-y_{2,g}z_{1,g^2}\equiv 0.$$ 

     We consider the endomorphism that takes the variables $y_{1,g^2}$ and $z_{1,g^2}$, respectively, into the $(G,*)$-polynomials $y_{1,g}\circ y_{3,g}$ and $[y_{1,g}, y_{3,g}]$, symmetric and skew, respectively, of homogeneous degree $g^2$. Adding up the identities we get 
$$y_{1,g}y_{3,g}y_{2,g} - y_{2,g}y_{1,g}y_{3,g}\equiv 0.$$

 Therefore, by Lemma \ref{ref0.3}, it follows that $m_{\langle \lambda \rangle}\le 1$, for all multipartition $\langle \lambda \rangle \vdash (n_{g^+})$ with $n\geq 3,$ concluding the proof of the lemma.
    \end{proof}

\end{document}
